\setcounter{reportfigure}{0}
\renewcommand{\thereportfigure}{C.\arabic{reportfigure}}
\section*{Appendix C --- Physical laboratory records}
Both runs used the laboratory model of Part~II (Appendix~B) with the course gains, which were not tuned.
The plots use the original 60 s trajectory: yaw above, pitch below. Yellow denotes the course reference and blue the measured response. The large excursions correspond to disturbances deliberately introduced by the operator's hands; their magnitude and duration were not measured.
\begin{center}
\includegraphics[width=\linewidth]{hardware_manual_disturbance_run_1.png}
\reportcaption{Physical run 1: course-reference tracking with manually introduced disturbances. Pronounced response excursions are visible around 20--25 s and 50--55 s. Original Scope screenshot.}
\end{center}
\begin{center}
\includegraphics[width=\linewidth]{hardware_manual_disturbance_run_2.png}
\reportcaption{Physical run 2: the same course trajectory with manual disturbances, including a pronounced response around 45--50 s. Original Scope screenshot.}
\end{center}
Both runs show qualitative recovery toward the moving reference after perturbation. Some pitch excursions exceed the displayed vertical range. Exact peak errors, disturbance-rejection settling times, RMS errors and voltage saturation cannot be established from these screenshots alone.
